% Capítulo 2: Marco Teórico
% Propósito: Presentar las bases teóricas, los conceptos clave y el contexto del sector de acopio pesquero que fundamentan el proyecto.
% Contenido: tomado del documento "TRABAJO DE CULMINACIÓN DE GRADO - AVANCES".

\section{Marco Teórico}

En el ámbito general internacionalmente, los centros de acopio de productos marinos constituyen la unidad operativa fundamental y el punto clave de la cadena de valor pesquera artesanal e intermedia, actuando como el puente principal entre las capturas desembarcadas en las costas y los mercados finales de consumo masivo y de exportación \parencite{fao2020,ifad2020,salas2011}. La dinámica diaria de trabajo en estos establecimientos inicia en las primeras horas de la jornada con la recepción de la materia prima en muelle, donde se realiza un proceso de selección y clasificación por especie, la categorización por calibres comercializables y una rigurosa evaluación del producto para determinar la frescura mediante el análisis de características críticas como brillo, textura de la carne, olor y firmeza agallada \parencite{garcia2021,sumaila2020}. De manera inmediata y como exigencia biosanitaria, el personal ejecuta la pesada e introduce los lotes en recipientes térmicos con agentes refrigerantes ---tales como termos isotérmicos con mezclas proporcionales de hielo y agua, o congeladores industriales--- con el propósito de garantizar la estricta continuidad de la cadena de frío, inhibir la proliferación bacteriana no deseada y mitigar el deterioro biológico acelerado inherente al producto biológico marino \parencite{fedepesca2020,shaw2019}.

Por lo tanto, en el contexto particular de países en desarrollo y concretamente en Nicaragua, la gestión operativa de los acopios pesqueros artesanales e intermedios, tanto en el Litoral del Pacífico como en las Regiones Autónomas de la Costa Caribe, mantiene un sesgo marcadamente analógico y empírico. El control del volumen de captura, la asignación y trazabilidad de lotes, así como el apartado de pedidos previa venta (preventas) y la distribución final a intermediarios o clientes locales, se administran de forma convencional mediante anotaciones manuales en pizarras externas, cuadernos escolares o libretas de papel \parencite{garcia2021,ifad2020,inpesca2022}. Este flujo de trabajo empírico carece de un sistema centralizado de información, lo que genera una alta vulnerabilidad ante desincronizaciones entre el inventario físico disponible e inventario registrado, pérdidas directas por inexactitudes en la medición del stock real, traspapelo de comandas de venta y distorsiones en los márgenes de ganancia operativa al no contar con un registro digitalizado e inmediato de los desembarques diarios \parencite{fao2020,sumaila2020}.

El Acopio Arrecife es un establecimiento acreditado y regulado por el Instituto Nicaragüense de la Pesca y Acuicultura (INPESCA) en San Juan del Sur, y ejerce un rol fundamental, pues es uno de los acopios de la cadena de suministro pesquera local mediante la adquisición diaria de especies marinas en el muelle pesquero, su posterior traslado al establecimiento y su conservación en termos con hielo y congeladores. No obstante, la operatividad cotidiana del acopio se sustenta en métodos puramente tradicionales y empíricos, como es el inventario físico, el cual se exhibe de forma estática en una pizarra acrílica, mientras que el registro contable de las transacciones se gestiona manualmente en cuadernos de apuntes. En la atención a sus clientes habituales (restaurantes, hoteles y familias), las reservas o preventas solicitadas por telefonía o mensajería (WhatsApp) no se formalizan en el inventario de manera inmediata, sino hasta el momento de la entrega física. Esta falta de sincronización en tiempo real genera desbalances en el control de existencias, desfasamiento en las cuentas y el riesgo operativo de sobreventa, limitando a su vez la generación de información estadística precisa sobre los volúmenes comercializados y la demanda por especies.

El fundamento teórico del presente proyecto se basa centralmente en la Teoría de Sistemas de Información Gerencial (SIG) propuesta por \textcite{laudon2020}, la cual establece que la automatización de la captura de datos en el punto de origen transforma la información operativa dispersa en un activo estratégico fundamental para la toma de decisiones. En el contexto del sector pesquero, este marco se complementa con los principios de la Trazabilidad en la Cadena de Suministro \parencite{gs1mexico2022}, donde se exige que cada lote de producto marino sea plenamente rastreable desde su ingreso físico a los recipientes de conservación hasta su venta final. Este principio sustenta la arquitectura del módulo de inventario por lotes adaptado para el Acopio «Arrecife».

En lo que respecta a los conceptos fundamentales que guían el desarrollo tecnológico del proyecto, un Sistema de Información Web se define como una aplicación bajo el modelo cliente-servidor cuya interfaz se ejecuta directamente dentro de un navegador mediante protocolos HTTP/HTTPS, mientras su lógica de negocio interactúa de forma remota con un servidor de base de datos relacional \parencite{sommerville2019}. Dentro de esta estructura, el Control de Inventario en Tiempo Real constituye un estado transaccional en el que cada evento de entrada, reserva o salida de mercancía actualiza de manera inmediata e indivisible el registro de existencias en el repositorio central, garantizando la consistencia permanente de las cifras disponibles \parencite{pressman2021}.

Por su parte, el concepto de Preventa o Apartado de Stock se entiende como un estado operativo temporal mediante el cual un volumen específico de producto se asigna de forma exclusiva a una orden pendiente de entrega, bloqueando su disponibilidad para otros compradores y previniendo la sobreventa \parencite{villasmil2021}. Para sostener esta lógica, la plataforma adopta la Clean Architecture (Arquitectura Limpia), un patrón de diseño de software introducido por Robert C. Martin que separa la lógica de negocio central de las interfaces de usuario, marcos de trabajo (frameworks) y motores de bases de datos, garantizando que el sistema sea mantenible y administrable de forma autónoma por el propio personal del acopio \parencite{martin2018}. Finalmente, la Persistencia Relacional asegura el almacenamiento estructurado de los datos mediante tablas vinculadas por claves primarias y foráneas, preservando la integridad referencial en cada transacción de compra y venta efectuada en el negocio \parencite{elmasri2020}.
